\documentclass[a4paper,12pt]{article}
\usepackage{amsmath, amsthm}
\usepackage{datetime}
\usepackage{framed}
\usepackage{enumitem}
\usepackage{fancyref}
\usepackage{wrapfig}
\usepackage{pifont}
\usepackage{appendix}
\usepackage{caption}
\usepackage{xcolor}
\usepackage[stable]{footmisc}
\usepackage{multicol}
\usepackage{csquotes}
\usepackage{pdfpages}

\usepackage{amsthm}

\usepackage{amssymb}
\usepackage{amsfonts}
\usepackage{amsmath}
\usepackage{mathtools}

\usepackage{tikz}
\usepackage{pgf}
\usepgflibrary{fpu}
\usepackage{qtree}
\usetikzlibrary{angles,fit,arrows,calc,math,intersections,through,backgrounds,cd}
\usepackage{pgfplots}
\usepackage{tkz-euclide}

\usepackage{listings}
\lstset{
  basicstyle=\itshape,
  xleftmargin=3em,
  literate={->}{$\rightarrow$}{2}
           {α}{$\alpha$}{1}
           {δ}{$\delta$}{1}
}


\usepackage{csquotes}
\renewcommand{\mkbegdispquote}[2]{\itshape}

\newdateformat{nianyueri}{修订于 \THEYEAR 年 \THEMONTH 月 \THEDAY 日 }

\usepackage{xstring}
\usepackage{catchfile}
\CatchFileDef{\HEAD}{../.git/refs/heads/master}{}
\newcommand{\gitrevision}{%
  \StrLeft{\HEAD}{7}%
}

\usepackage{data/quiver}
\usepackage{data/circledsteps}
\usepackage[top=1in,bottom=1in,left=1in,right=1in]{geometry} % 用于设置页面布局
\usepackage{xeCJK} % 用于使用本地字体
\usepackage[super, square, sort&compress]{natbib} % 处理参考文献
\usepackage{titlesec, titletoc} % 设置章节标题及页眉页脚
%\usepackage{xCJKnumb} % 中英文数字转换
\usepackage{amssymb}
\usepackage{amsmath} % 在公式中用\text{文本}输入中文
\usepackage{diagbox}
\usepackage{multirow} % 表格中使用多行
\usepackage{booktabs} % 表格中使用\toprule等命令
\usepackage{rotating} % 使用sidewaystable环境旋转表格
\usepackage{tabularx}
\usepackage{graphicx} % 处理图片
\usepackage{footnote} % 增强的脚注功能，可添加表格脚注
\usepackage{threeparttable} % 添加真正的表格脚注,示例见README
\usepackage{hyperref} % 添加pdf书签

\usepackage{tikz}
\usetikzlibrary{shapes,arrows,shadows}

% 字体设置
\setmainfont{Times New Roman}
\setsansfont[Scale=MatchLowercase,Mapping=tex-text]{PT Sans}
\setmonofont[Scale=MatchLowercase]{PT Mono}
\setCJKmainfont[ItalicFont={Kaiti SC}, BoldFont={Heiti SC}]{Songti SC}
\setCJKsansfont{Heiti SC}
\setCJKmonofont{Songti SC}
% \setCJKmainfont[BoldFont={FZXiaoBiaoSong-B05S}]{Songti SC}
% \setCJKfamilyfont{kai}[BoldFont=Heiti SC]{Kaiti SC}
% \setCJKfamilyfont{song}[BoldFont=Heiti SC]{Songti SC}
% \setCJKfamilyfont{hei}[BoldFont=Heiti SC]{Heiti SC}
% \setCJKfamilyfont{fsong}[BoldFont=Heiti SC]{Songti SC}
% \newcommand{\kai}[1]{{\CJKfamily{kai}#1}}
% \newcommand{\hei}[1]{{\CJKfamily{hei}#1}}
% \setromanfont[Mapping=tex-text]{TeXGyrePagella}
% \setsansfont[Scale=MatchLowercase,Mapping=tex-text]{TeXGyrePagella}
% \setmonofont[Scale=MatchLowercase]{Courier New}
%%设置常用中文字号，方便调用
\newcommand{\erhao}{\fontsize{22pt}{\baselineskip}\selectfont}
\newcommand{\xiaoerhao}{\fontsize{18pt}{\baselineskip}\selectfont}
\newcommand{\sanhao}{\fontsize{16pt}{\baselineskip}\selectfont}
\newcommand{\xiaosanhao}{\fontsize{15pt}{\baselineskip}\selectfont}
\newcommand{\sihao}{\fontsize{14pt}{\baselineskip}\selectfont}
\newcommand{\xiaosihao}{\fontsize{12pt}{\baselineskip}\selectfont}
\newcommand{\wuhao}{\fontsize{10.5pt}{\baselineskip}\selectfont}
\newcommand{\xiaowuhao}{\fontsize{9pt}{\baselineskip}\selectfont}
\newcommand{\liuhao}{\fontsize{7.5pt}{\baselineskip}\selectfont}

% 章节标题显示方式及页眉页脚设置
% \item xCJKnumb是自己额外安装的包
% \item titleformat命令定义标题的形式
% \item titlespacing定义标题距左、上、下的距离
\titleformat{\section}{\raggedright\large\bfseries}{\thesection}{1em}{}
\titleformat{\subsection}{\raggedright\normalsize\bfseries}{\thesubsection}{1em}{}
\titlespacing{\section}{0pt}{*0}{*2}
\titlespacing{\subsection}{0pt}{*0}{*1}
% 由于默认的2em缩进不够，所以我手动调整了，但是在windows下似乎2.2就差不多了，或者是article中没有这个问题
\setlength{\parindent}{2.2em}

% 设置表格标题前后间距
\setlength{\abovecaptionskip}{0pt}
\setlength{\belowcaptionskip}{0pt}


\renewcommand{\refname}{\bfseries{参~考~文~献}} %将Reference改为参考文献（用于 article）
% \renewcommand{\bibname}{参~考~文~献} %将bibiography改为参考文献（用于 book）
\renewcommand{\baselinestretch}{1.38} %设置行间距
\renewcommand{\figurename}{\small\ttfamily 图}
\renewcommand{\tablename}{\small\ttfamily 表}


\usepackage{stmaryrd}
\usepackage{mathtools}
\usepackage{wasysym}
\usepackage{textcomp}
\usepackage{blindtext}
\usepackage{subfiles}

\newtheorem{problem}{问题}
\numberwithin{problem}{section}
\newtheorem{definition}{定义}
\numberwithin{definition}{section}
\newtheorem{lemma}{引理}
\numberwithin{lemma}{section}
\newtheorem{proposition}{命题}
\numberwithin{proposition}{section}
\newtheorem{theorem}{定理}
\numberwithin{theorem}{section}
\newtheorem{grammar}{文法}
\numberwithin{grammar}{section}
\newtheorem{program}{程序}
\numberwithin{program}{section}
\newtheorem{convention}{约定}
\numberwithin{convention}{section}
\newtheorem{corollary}{推论}
\numberwithin{corollary}{section}
\renewcommand*{\proofname}{证明}

\xeCJKsetwidth{‘’“”}{1em}

\title{我与集智}
\date{\nianyueri\today}
\author{苑明理}

\begin{document}

\begingroup
\let\newpage\relax
\maketitle
\endgroup

\centerline{\rule{13cm}{0.4pt}}
\renewcommand{\contentsname}{\hfill\bfseries 目录\hfill}
\setcounter{tocdepth}{2}
\tableofcontents
\centerline{\rule{13cm}{0.4pt}}

\section{初识集智}

早就在网上听说过集智俱乐部的大名了，2012 年我再次回到北京，其后不久，就和吴令飞、袁行远约好在对外经贸大学外的咖啡馆碰碰头。
还记得当时吴令飞讲了渗流中的重整化、注意力流，我则说了我理解的维基第一原则、概念从模糊到清晰的颗粒化过程、瓦克星计划，
然后商定了我主讲一次可计算性理论，按照集智的加入规则，就可以成为集智的成员了。

\section{古北水镇：第一次研读营}

\begin{displayquote}
“十余年程序员生涯，游走于工程、语言、知识、智能、数学这些关键词之间，内心对世界充满很多好奇和困惑，尝听友人讲大刘《山》的故事，
也期自己可凭蛮力，凿空厚壁，得见星空。”
\end{displayquote}

\section{古北水镇：第二次研读营}

\section{喜洲：第三次研读营}

\section{线上：第四次研读营}

机器能否数数的思考

2021 年 1 月集智年会上我做了一个小的分享，去探讨机器能不能数数的问题。

表面上看似乎不是问题，机器本来不就是高级版本的计算器嘛！实则不然，我其实在这里提了
一个不同于图灵测试的智能体的创造性测试问题。图灵测试是以人类智能为标准的，而这里的
机器数数的测试问题则是以一种超出了人类智能的方式来设问的。

在演示文档中，我有一层层的剥开这个问题，先考察生物和认知里的数，继而从目前人类的形
式定义去考察，发现想把这个问题定义清楚都比较困难。 困难在于，无法简单的排除掉智能体
可以创造出和人类不同的数的形式，但这种情况下，你无法给出一个数的标准定义，从而也很
难给出能不能数数问题的测试标准。 最后一段，我尝试用强化学习的思路，创造一个生存环境，
智能体必须通过计数的方法来获得更优的选择，然后不限定智能体的策略的制定，如此从侧面
去考察数的创造。

这种强化学习的思路，最终还是通过报酬或者激励的量化来驱动策略的搜索，也就是把数的创造
问题化归为一个势能面上的寻优。这背后是一种\textbf{几何式的元数学}观点。

扩展开来，这里面其实有个很深的智能体之间彼此理解的问题。
我对沙漠蚁计步的能力就很好奇：沙漠蚁在觅食时可以走很复杂的路线，但是它们一旦归巢，
就会直奔巢穴的位置。 沙漠蚁在这里展现出来的空间路径积分的能力，其实不比人类最开始掰
手指头数那么几百个数要低。我想的是数学能力其实是一种具身性的能力，而数学形式也受具
身性的影响。 想要隔着物种或者其他的隔阂去理解的话，必须得深入到那个具身环境里去理
解。但我倒不是否定柏拉图理念式的数学，我是想探索这些不同具身环境下的数学形态的某种
公约数。当生物逐渐从昆虫演化到了哺乳动物，它们在海马体里出现了导航神经细胞，这些细
胞有六边形的活跃模式。那么这种神经层面上六边形活跃模式对应的几何， 怎么和我们心理
层面的几何挂上钩的呢？从数学角度如何理解呢？所以还可以跨越神经、认知和数学，去思考
更多。

上面提到的公约数可能也是理解人类符号化知识的一个有趣切入的方向。考察这种异与同背后，
其实还有一个问题是什么是真实？有没有这种真实？ 如果没有这种真实，智能体之间能互相理
解吗？这背后的指涉，远不止在具身性数学的探讨，如今混乱的人类世界或许也需要深刻反思
这个问题。

\section{线上：第五次研读营}

\begin{figure}[ht]
\centering
\begin{tikzpicture}
\draw (6, 2) node[inner sep=0] {\includegraphics[width=4in]{images/ennepers.png}};
\node [anchor=225, black] at (9,6) {几何};
\node [anchor=225, black] at (9,-3) {分析};
\node [anchor=225, black] at (2,6) {计算};
\node [anchor=225, black] at (0,0) {学习};
\node [anchor=225, black] at (6,2) {算术表达式几何};
\end{tikzpicture}
\caption{知识的极小曲面}
\end{figure}

康德有论述“先天综合判断”，几何作为命题都呈现为“综合命题”，肥皂泡的几何式知识观有可能会暗合康德的理论。
这个知识观可以形式化为上面说的“过程与形式”的数学，而算术表达式几何就是这个“过程与形式”的数学的第一个例子。

\end{document}